% Problem 10 — Preconditioned Conjugate Gradient for the RKHS-constrained CP subproblem
% Solution by Claude (Opus 4.7), 2026-04-30
\documentclass[12pt]{article}
\usepackage{fullpage,amsmath,amssymb,amsthm,hyperref}
\usepackage{algorithm,algpseudocode}
\newtheorem{theorem}{Theorem}
\newtheorem{lemma}[theorem]{Lemma}
\newtheorem{proposition}[theorem]{Proposition}
\theoremstyle{definition}
\newtheorem{definition}[theorem]{Definition}
\newtheorem{remark}[theorem]{Remark}

\newcommand{\vecop}{\operatorname{vec}}
\newcommand{\matop}{\operatorname{mat}}
\newcommand{\trace}{\operatorname{tr}}
\newcommand{\range}{\operatorname{range}}
\newcommand{\diag}{\operatorname{diag}}
\newcommand{\op}{\mathrm{op}}
\newcommand{\F}{\mathrm{F}}
\newcommand{\R}{\mathbb{R}}
\newcommand{\Ht}{\widetilde{H}}
\newcommand{\Kt}{\widetilde{K}}
\newcommand{\Zt}{\widetilde{Z}}

\title{Solution to Problem 10: Preconditioned Conjugate Gradient \\
       for the RKHS-Constrained CP Subproblem}
\author{}
\date{}

\begin{document}
\maketitle

\section*{Setup and notation}

We follow the problem statement of Kolda--Ward. Let $\mathcal{T}\in\R^{n_1\times\cdots\times n_d}$
be a $d$-way tensor with missing entries. Fix a mode $k$, set
\[
n \;:=\; n_k,\qquad
M \;:=\; \prod_{i\neq k} n_i,\qquad
N \;:=\; \prod_{i=1}^d n_i \;=\; nM,
\]
and let $q$ denote the number of \emph{observed} entries. We assume
\begin{equation}\label{eq:size-assump}
d \;<\; n,\quad r \;<\; q,\qquad q \;\ll\; N,\qquad n \;\ll\; M.
\end{equation}
The selection matrix $S\in\R^{N\times q}$ is a column-subset of $I_N$: $S^T x\in\R^q$ extracts
the $q$ observed entries. Let $Z\in\R^{M\times r}$ be the Khatri--Rao product of all factor
matrices except mode $k$ (that is, $Z=A_d\odot\cdots\odot A_{k+1}\odot A_{k-1}\odot\cdots\odot A_1$),
let $K\in\R^{n\times n}$ be a positive semidefinite RKHS kernel matrix, and let $\lambda>0$ be the
ridge parameter. Writing the unknown factor as $A_k=KW$ with $W\in\R^{n\times r}$, the
mode-$k$ subproblem is the linear system
\begin{equation}\label{eq:main-system}
\mathbf{H}\,\vecop(W) \;=\; (I_r\otimes K)\,\vecop(B),\qquad
B \;=\; T Z,
\end{equation}
where $T\in\R^{n\times M}$ is the mode-$k$ unfolding (with missing entries set to $0$), and
\begin{equation}\label{eq:H-def}
\mathbf{H} \;:=\; (Z\otimes K)^T\,SS^T\,(Z\otimes K) \;+\; \lambda\,(I_r\otimes K).
\end{equation}
The system has size $nr\times nr$. Throughout we write $v=\vecop(V)$ with $V\in\R^{n\times r}$
(column-major: $\vecop(V)$ stacks the columns of $V$).
We use the Kronecker--vec identity
\begin{equation}\label{eq:kron-vec}
(B^T\otimes A)\,\vecop(V) \;=\; \vecop(A V B).
\end{equation}

\paragraph{Index conventions.}
The selection matrix $S$ encodes a set $\Omega\subseteq[n]\times[M]$ with $|\Omega|=q$.
Each observation $\omega\in\Omega$ corresponds to a pair $(i_\omega,j_\omega)$ with
$i_\omega\in[n]$, $j_\omega\in[M]$. By column-major vectorization, the $\omega$-th row of
$S^T$ selects index $i_\omega+(j_\omega-1)n$ from $\vecop(\cdot)$. We let $\Kt\in\R^{q\times n}$
denote the $q$ rows of $K$ indexed by $\{i_\omega:\omega\in\Omega\}$ (with multiplicity, in the
order of $\Omega$), and $\Zt\in\R^{q\times r}$ the $q$ rows of $Z$ indexed by
$\{j_\omega:\omega\in\Omega\}$ (also with multiplicity, in the same order).

\paragraph{Outline.}
Section~\ref{sec:proper} establishes that $\mathbf{H}\succ 0$, hence CG applies and is
well-defined.
Section~\ref{sec:direct} states the direct symmetric (Cholesky) method as Algorithm~\ref{alg:direct}.
Section~\ref{sec:matvec} proves the key matrix-vector product proposition: $\mathbf{H}v$ can be
formed in $O(n^2r+qr)$ using only row-subset accesses $\Kt,\Zt$, never any $O(M)$ or $O(N)$
operation.
Section~\ref{sec:precond} defines the preconditioner $P=\lambda(I_r\otimes K)$ and bounds the
condition number $\kappa(P^{-1}\mathbf{H})$.
Section~\ref{sec:pcg} presents PCG as Algorithm~\ref{alg:pcg}.
Section~\ref{sec:cost} compiles the comparative complexity statement.

\section{Properties of the system}\label{sec:proper}

\begin{proposition}[Symmetry and positive definiteness]\label{prop:psd}
Assume $K\succ 0$ and $\lambda>0$. Then $\mathbf{H}$ is symmetric positive definite.
\end{proposition}

\begin{proof}
\emph{Symmetry.} Both summands of \eqref{eq:H-def} are symmetric:
$(Z\otimes K)^T(SS^T)(Z\otimes K)$ is the Gram of the matrix $S^T(Z\otimes K)$, hence symmetric
PSD; and $\lambda(I_r\otimes K)$ is symmetric since $K$ is.

\emph{Positive definiteness.} Take any $v\in\R^{nr}\setminus\{0\}$. Then
\[
v^T \mathbf{H} v
\;=\; \|S^T(Z\otimes K)v\|_2^2 \;+\; \lambda\,v^T(I_r\otimes K)v.
\]
The first term is $\geq 0$. For the second: $v^T(I_r\otimes K)v=\sum_{j=1}^r v_j^T K v_j$ with
$v_j\in\R^n$ the $j$-th block of $v$. Since $K\succ 0$, each $v_j^T K v_j\geq 0$ and
$v_j^T K v_j=0$ iff $v_j=0$. Because $v\neq 0$ at least one $v_j\neq 0$, so
$v^T(I_r\otimes K)v>0$, and thus $v^T\mathbf{H}v>\lambda\cdot v^T(I_r\otimes K)v/2>0$
when $\lambda>0$. Hence $\mathbf{H}\succ 0$.
\end{proof}

\begin{remark}\label{rmk:psd-only}
If $K$ is only PSD with nullspace $\mathcal{N}\neq\{0\}$, then $\mathbf{H}$ is still PSD and is
positive definite on the orthogonal complement of $\mathcal{N}^{\otimes r}$; the right-hand side
$(I_r\otimes K)\vecop(B)$ lies in $\range(I_r\otimes K)$, so the system is consistent and the
unique minimum-norm solution is well-defined. CG then converges on the active subspace. For
notational simplicity we work under $K\succ 0$ in what follows; the rank-deficient case requires
only replacing $K^{-1}$ by the Moore--Penrose pseudoinverse $K^\dagger$ in the preconditioner.
\end{remark}

By Proposition~\ref{prop:psd} the symmetric direct method (Cholesky factorization) and the
(preconditioned) conjugate gradient method are both applicable.

\section{Direct symmetric method}\label{sec:direct}

\subsection{Construction of $\mathbf{H}$ via row subsets}

We first record a compact formula for $\mathbf{H}$ that uses only $\Kt$ and $\Zt$.

\begin{lemma}[Row-subset form of $\mathbf{H}$]\label{lem:rowsubset}
With the notation above,
\begin{equation}\label{eq:rowsub}
(Z\otimes K)^T SS^T (Z\otimes K) \;=\; \sum_{\omega\in\Omega}
\bigl(\,\Zt(\omega,:)^T \otimes \Kt(\omega,:)^T\,\bigr)\bigl(\,\Zt(\omega,:)\otimes \Kt(\omega,:)\,\bigr).
\end{equation}
Equivalently, the $((i,s),(i',s'))$-block of size $1\times 1$ of $(Z\otimes K)^TSS^T(Z\otimes K)$
is
\begin{equation}\label{eq:Hentry}
\mathbf{H}^{(0)}_{(i,s),(i',s')} \;:=\; \sum_{\omega\in\Omega} \Kt(\omega,i)\,\Zt(\omega,s)\,\Kt(\omega,i')\,\Zt(\omega,s')
\end{equation}
for $i,i'\in[n]$, $s,s'\in[r]$, where $\mathbf{H}^{(0)}$ is the first summand of \eqref{eq:H-def}.
\end{lemma}

\begin{proof}
The matrix $S^T$ extracts the $q$ observed entries; the $\omega$-th row of $S^T$ is the standard
basis row indexed by $i_\omega+(j_\omega-1)n$. Therefore
\[
S^T(Z\otimes K)
\;=\;
\bigl[\,(Z\otimes K)_{\,i_\omega+(j_\omega-1)n,\,\cdot}\,\bigr]_{\omega\in\Omega}.
\]
By the Kronecker definition, $(Z\otimes K)_{i+(j-1)n,\,a+(s-1)n}=Z_{j,s}K_{i,a}$. The
$\omega$-th row of $S^T(Z\otimes K)$ is therefore
$\Zt(\omega,:)\otimes \Kt(\omega,:)$, an $1\times nr$ row vector. Multiplying by its transpose
on the left and summing over $\omega$ gives \eqref{eq:rowsub}. The entrywise form
\eqref{eq:Hentry} follows by reading off the $((i,s),(i',s'))$ entry of each rank-one summand.
\end{proof}

\begin{proposition}[Cost of forming $\mathbf{H}$]\label{prop:formH}
Forming the symmetric matrix $\mathbf{H}\in\R^{nr\times nr}$ explicitly costs
$O(qn^2r^2)$ time and $O(n^2r^2)$ memory.
\end{proposition}

\begin{proof}
By \eqref{eq:Hentry}, each of the $\binom{nr+1}{2}\le n^2r^2$ upper-triangular entries of
$\mathbf{H}^{(0)}$ is a sum over $\omega\in\Omega$ of a product of four scalars. Computing one
entry costs $O(q)$, so all entries cost $O(qn^2r^2)$. Adding the diagonal-block correction
$\lambda(I_r\otimes K)$ adds $rn^2$ entries each computed in $O(1)$, an $O(n^2r)$ contribution
that is dominated by $qn^2r^2$ since $r\geq 1$ and $q\geq 1$.
The matrix $\mathbf{H}$ has $n^2r^2$ entries (after symmetrization), giving the stated memory.
\end{proof}

\subsection{Algorithm box: direct symmetric method}

\begin{algorithm}[ht]
\caption{Direct symmetric solver for $\mathbf{H}\,\vecop(W)=(I_r\otimes K)\vecop(B)$}\label{alg:direct}
\begin{algorithmic}[1]
\Require Kernel $K\in\R^{n\times n}$ with $K\succ 0$; Khatri--Rao $Z\in\R^{M\times r}$;
selection set $\Omega=\{\omega=(i_\omega,j_\omega)\}_{\omega=1}^q$; observed values
$y\in\R^q$ with $y_\omega=T_{i_\omega,j_\omega}$ (the observed tensor entries);
ridge $\lambda>0$.
\Ensure Solution $W\in\R^{n\times r}$ to \eqref{eq:main-system}.
\State \textbf{Row-subset extraction.}\;
       Form $\Kt\in\R^{q\times n}$ with $\Kt(\omega,:)=K(i_\omega,:)$ and
       $\Zt\in\R^{q\times r}$ with $\Zt(\omega,:)=Z(j_\omega,:)$.
       \Comment{cost $O(qn+qr)$}
\State \textbf{Right-hand side.}\;
       Compute $B\in\R^{n\times r}$ via $B(i,s)=\sum_{\omega:i_\omega=i} y_\omega \Zt(\omega,s)$,
       i.e.\ scatter-add the contributions of observations into mode-$k$ slots.
       \Comment{cost $O(qr)$}
\State Form $C\leftarrow K\,B \in\R^{n\times r}$, then $g \leftarrow \vecop(C)\in\R^{nr}$.
       \Comment{$g=(I_r\otimes K)\vecop(B)$; cost $O(n^2 r)$}
\State \textbf{Build $\mathbf{H}$.}\;
       Initialize $\mathbf{H}\leftarrow 0\in\R^{nr\times nr}$.
\For{$\omega = 1$ \textbf{to} $q$}
   \State $\kappa_\omega \leftarrow \Kt(\omega,:)\in\R^{n}$;\quad
          $\zeta_\omega \leftarrow \Zt(\omega,:)\in\R^{r}$.
   \State Update $\mathbf{H}\mathrel{+}= (\zeta_\omega\zeta_\omega^T)\otimes(\kappa_\omega\kappa_\omega^T)$.
          \Comment{rank-one block update; cost $O(n^2r^2)$}
\EndFor
\State Add ridge: $\mathbf{H}\mathrel{+}=\lambda\,(I_r\otimes K)$.\quad
       \Comment{loop above costs $O(qn^2r^2)$ (Proposition~\ref{prop:formH}); ridge add $O(n^2 r)$}
\State \textbf{Cholesky factorize}: $\mathbf{H} = L L^T$ with $L\in\R^{nr\times nr}$ lower
       triangular.\quad \Comment{cost $O(n^3 r^3)$}
\State Solve $L u = g$ by forward substitution.\quad \Comment{cost $O(n^2 r^2)$}
\State Solve $L^T w = u$ by back substitution.\quad \Comment{cost $O(n^2 r^2)$}
\State \Return $W \leftarrow \matop(w)\in\R^{n\times r}$ where $w=\vecop(W)$.
\end{algorithmic}
\end{algorithm}

\begin{proposition}[Correctness, termination, and complexity of Algorithm~\ref{alg:direct}]
\label{prop:direct-cost}
Algorithm~\ref{alg:direct} terminates and returns the unique $W$ solving \eqref{eq:main-system}.
Its time complexity is $O(qn^2r^2+n^3r^3)$ and its memory is $O(n^2r^2)$.
\end{proposition}

\begin{proof}
\emph{Correctness.} Step~3 computes $g=(I_r\otimes K)\vecop(B)$ since by \eqref{eq:kron-vec},
$\vecop(K B I_r)=(I_r\otimes K)\vecop(B)$. The loop in Steps 5--7 implements the rank-one
expansion of Lemma~\ref{lem:rowsubset}: each iteration adds
$(\zeta_\omega\zeta_\omega^T)\otimes(\kappa_\omega\kappa_\omega^T)
=(\zeta_\omega\otimes\kappa_\omega)(\zeta_\omega\otimes\kappa_\omega)^T$, the rank-one matrix
appearing in \eqref{eq:rowsub}. Adding $\lambda(I_r\otimes K)$ yields exactly $\mathbf{H}$ from
\eqref{eq:H-def}. By Proposition~\ref{prop:psd}, $\mathbf{H}\succ 0$ so Cholesky succeeds; the
forward and back substitutions then return $w=\mathbf{H}^{-1}g$, the unique solution.

\emph{Termination.} All loops are over finite index sets. Cholesky on a symmetric positive
definite matrix terminates in finitely many steps.

\emph{Complexity.} The for-loop performs $q$ rank-one outer-product updates of an
$nr\times nr$ matrix, each $O(n^2r^2)$, totaling $O(qn^2r^2)$. Cholesky on an $nr\times nr$
matrix is $O((nr)^3)=O(n^3r^3)$. Forward/back substitution is $O((nr)^2)$. Steps 1--3 are
linear in the inputs ($O(qn+qr+n^2r)$ which is dominated). The total is
$O(qn+qr+n^2r+qn^2r^2+n^3r^3)=O(qn^2r^2+n^3r^3)$ under \eqref{eq:size-assump}.
Memory is dominated by storing $\mathbf{H}$ and $L$: $2\cdot n^2r^2$ floats, i.e.\ $O(n^2r^2)$.
\end{proof}

\section{Matrix--vector product without $O(M)$ or $O(N)$ work}\label{sec:matvec}

The PCG method only requires the action $v\mapsto \mathbf{H}v$, never the explicit matrix.
The next proposition is the workhorse: each application costs $O(n^2 r+qr)$ time and accesses
$Z$ and $K$ only through the row-subset matrices $\Zt$ and $\Kt$.

\begin{proposition}[Fast matrix-vector product]\label{prop:matvec}
Let $v=\vecop(V)$ with $V\in\R^{n\times r}$. Let $\Kt\in\R^{q\times n}$ and
$\Zt\in\R^{q\times r}$ be the row-subset matrices defined above. Then
\begin{equation}\label{eq:Hv-formula}
\matop(\mathbf{H}v) \;=\; K\bigl(G(V) + \lambda V\bigr) \;\in\;\R^{n\times r},
\end{equation}
where $G(V)\in\R^{n\times r}$ is computed by:
\begin{enumerate}
\item[(A)] $\Ht\leftarrow \Kt\, V \in\R^{q\times r}$ (so that $\Ht(\omega,:)=\Kt(\omega,:)V=
(KV)(i_\omega,:)$);
\item[(B)] $u\in\R^q$ with $u_\omega = \langle \Ht(\omega,:),\,\Zt(\omega,:)\rangle$ for each
$\omega\in\Omega$;
\item[(C)] $G(V)\in\R^{n\times r}$ with $G(V)(i,:)=\sum_{\omega:\,i_\omega=i} u_\omega\,\Zt(\omega,:)$.
\end{enumerate}
The total cost of computing \eqref{eq:Hv-formula} from $V$ is $O(n^2 r+qr)$, and only $\Kt,\Zt$
are accessed (no $O(M)$ or $O(N)$ work).
\end{proposition}

\begin{proof}
\emph{Step (i): Ridge term.}
By \eqref{eq:kron-vec}, $\lambda(I_r\otimes K)v=\lambda(I_r\otimes K)\vecop(V)
=\lambda\,\vecop(K V I_r)=\vecop(\lambda K V)$, so $\matop$ of the ridge contribution is
$\lambda KV$. We will absorb this into the final left multiplication by $K$.

\emph{Step (ii): Data term.}
By Lemma~\ref{lem:rowsubset},
\[
(Z\otimes K)^T SS^T (Z\otimes K)\,v
\;=\; \sum_{\omega\in\Omega} \bigl(\Zt(\omega,:)^T\otimes \Kt(\omega,:)^T\bigr)\,
\bigl(\Zt(\omega,:)\otimes \Kt(\omega,:)\bigr)\,v.
\]
By \eqref{eq:kron-vec}, $(\Zt(\omega,:)\otimes \Kt(\omega,:))\vecop(V)
=\vecop\bigl(\Kt(\omega,:)\,V\,\Zt(\omega,:)^T\bigr)
=\Kt(\omega,:)\,V\,\Zt(\omega,:)^T$, a scalar (since $\Kt(\omega,:)\in\R^{1\times n}$,
$V\in\R^{n\times r}$, $\Zt(\omega,:)^T\in\R^{r\times 1}$). Calling this scalar $u_\omega$,
\begin{align}
u_\omega
&= \Kt(\omega,:)\,V\,\Zt(\omega,:)^T \notag\\
&= \bigl(\Kt(\omega,:)\,V\bigr)\,\Zt(\omega,:)^T
\;=\; \langle \Ht(\omega,:),\,\Zt(\omega,:)\rangle, \label{eq:u-def}
\end{align}
which is exactly Step~(A) followed by Step~(B): first form $\Ht=\Kt V$ row-by-row, then
inner-product each row against the corresponding row of $\Zt$.

Now applying $(\Zt(\omega,:)^T\otimes \Kt(\omega,:)^T)$ to the scalar $u_\omega$ gives
$u_\omega\bigl(\Zt(\omega,:)^T\otimes\Kt(\omega,:)^T\bigr)\in\R^{nr}$, which by
\eqref{eq:kron-vec} equals $\vecop\bigl(u_\omega\,\Kt(\omega,:)^T\,\Zt(\omega,:)\bigr)$.
Therefore
\begin{equation}\label{eq:matop-data}
\matop\Bigl[(Z\otimes K)^T SS^T (Z\otimes K)\,v\Bigr]
\;=\; \sum_{\omega\in\Omega} u_\omega\,\Kt(\omega,:)^T\,\Zt(\omega,:).
\end{equation}
The matrix $\Kt(\omega,:)^T\,\Zt(\omega,:)\in\R^{n\times r}$ has $i$-th row equal to
$\Kt(\omega,i)\,\Zt(\omega,:)$. The $i$-th row of \eqref{eq:matop-data} is therefore
\[
\sum_{\omega\in\Omega} u_\omega\,\Kt(\omega,i)\,\Zt(\omega,:)
\;=\; \sum_{\omega\in\Omega} u_\omega\,K(i_\omega,i)\,\Zt(\omega,:).
\]
Equivalently, defining
$\widehat{G}(i_\omega,:):=\sum_{\omega'\in\Omega:\,i_{\omega'}=i_\omega} u_{\omega'}\Zt(\omega',:)$
as the scatter-sum onto the row index $i_\omega\in[n]$, and recognizing that
$\Kt(\omega,i)=K(i_\omega,i)$, we can rewrite \eqref{eq:matop-data} as $K\,G(V)$, where
\begin{equation}\label{eq:G-def}
G(V)(i,:) \;:=\; \sum_{\omega\in\Omega:\,i_\omega=i} u_\omega\,\Zt(\omega,:),\qquad i=1,\dots,n.
\end{equation}
Indeed the $a$-th row of $K\,G(V)$ is $\sum_{i=1}^n K(a,i)\,G(V)(i,:)
=\sum_{i=1}^n\sum_{\omega:i_\omega=i}K(a,i)u_\omega\Zt(\omega,:)
=\sum_{\omega\in\Omega}K(a,i_\omega)u_\omega\Zt(\omega,:)$, and $K(a,i_\omega)=\Kt(\omega,a)$,
matching \eqref{eq:matop-data} row-by-row. This is Step~(C).

\emph{Step (iii): Combine.}
Adding the ridge term $\lambda KV$ from Step~(i) yields
\[
\matop(\mathbf{H}v) \;=\; K\,G(V) + \lambda K V \;=\; K\bigl(G(V)+\lambda V\bigr),
\]
which is \eqref{eq:Hv-formula}.

\emph{Cost.}
\begin{itemize}
\item Step~(A) ($\Ht=\Kt V$): the matrix $\Kt$ is $q\times n$, $V$ is $n\times r$, so the product
costs $O(qnr)$. We never form $KV$ entirely; we only need its rows indexed by $\Omega$.
However we may also compute $KV$ in full (cost $O(n^2 r)$) and then read off
$\Ht(\omega,:)=(KV)(i_\omega,:)$. Both options are admissible. In our complexity bound we use
$O(n^2 r)$ for this step (since we will need $KV$ inside $K(G(V)+\lambda V)=KG(V)+\lambda KV$
in Step~(iii) anyway; we save a multiplication by reusing it). Alternatively, computing $\Ht$
directly via $\Kt V$ costs $O(qnr)$, which under the assumption $q\geq r$ in
\eqref{eq:size-assump} satisfies $qnr\geq nr^2$ but is unrelated to $n^2r$.
\item Step~(B) ($u\in\R^q$): for each of $q$ values of $\omega$, an inner product of two
$r$-vectors. Cost $O(qr)$.
\item Step~(C) ($G(V)\in\R^{n\times r}$): for each $\omega\in\Omega$ we add the scaled row
$u_\omega\Zt(\omega,:)$ to row $i_\omega$ of $G(V)$. This is one scatter-add of $r$ scalars per
$\omega$, total $O(qr)$.
\item Final left-multiplication by $K$ (Step~(iii)): $K(G(V)+\lambda V)\in\R^{n\times r}$ costs
$O(n^2 r)$. (Equivalently, compute $K G(V)$ in $O(n^2 r)$ and reuse the previously formed $KV$
to obtain $\lambda KV$ at no extra dominant cost.)
\end{itemize}
Total: $O(n^2 r) + O(qnr) + O(qr) + O(n^2 r) = O(n^2 r + qnr)$. Under
\eqref{eq:size-assump}, $qnr$ may dominate $qr$; if we instead use the alternative where
Step~(A) is computed via $\Kt V$ in $O(qnr)$ and we form $KV$ once globally in $O(n^2r)$ for
later reuse, the total is still $O(n^2r + qnr)$. Since the final answer involves
$\Ht(\omega,:)V$ as an $r$-vector for each $\omega$ (touching $n$ entries of $V$ per row), at
least $\Theta(qnr)$ work is unavoidable when the rows $\{i_\omega\}_{\omega\in\Omega}$ are
distinct; but in the typical regime where one precomputes the full matrix $KV$ once
($O(n^2r)$, independent of $q$) and then reads off rows $\Ht(\omega,:)=(KV)(i_\omega,:)$ at
$O(qr)$ total, the cost reduces to
\begin{equation}\label{eq:matvec-cost}
\boxed{\;O(n^2 r + qr)\;}
\end{equation}
as claimed in the proposition. We adopt this implementation in Algorithm~\ref{alg:pcg} below.

\emph{Memory access.}
The only places where $K$ or $Z$ enter are: (i) the full multiplication $KV$, which uses
$K\in\R^{n\times n}$ but never any column or row of $Z$; (ii) the row reads
$\Zt(\omega,:)=Z(j_\omega,:)$ used in Steps~(B) and (C). Both $\Kt$ and $\Zt$ have $q$ rows.
Nowhere does the procedure touch any other rows of $Z$ or any entries of $S$, $T$ outside of
$\Omega$. In particular, no operation has cost depending on $M$ or on $N$.
\end{proof}

\begin{remark}\label{rmk:matvec-savings}
The crucial step is precomputing $KV$ (cost $O(n^2r)$, \emph{independent of $q$}) and then
reading off only the $q$ rows $(KV)(i_\omega,:)$ via the row-subset matrix $\Kt$. Because
$n\ll M$ and $q\ll N=nM$, we have $n^2r$ much smaller than $Mr$ and $qr$ much smaller than $Nr$:
the algorithm avoids the prohibitive $O(M r)$ cost of forming the full Khatri--Rao product or
the $O(N)$ cost of touching every tensor entry.
\end{remark}

\section{Preconditioner and condition-number bound}\label{sec:precond}

\subsection{Choice of preconditioner}

The natural choice is
\begin{equation}\label{eq:P-def}
P \;:=\; \lambda\,(I_r\otimes K)
\end{equation}
which is the second (ridge) summand of $\mathbf{H}$. It is symmetric positive definite under
$K\succ 0$, $\lambda>0$, and its inverse acts blockwise:
\begin{equation}\label{eq:Pinv}
P^{-1}\,\vecop(Z) \;=\; \frac{1}{\lambda}\,\vecop\!\bigl(K^{-1} Z\bigr),
\qquad Z\in\R^{n\times r}.
\end{equation}

\begin{proposition}[Cost of one preconditioner solve]\label{prop:Pinv-cost}
Given the Cholesky factorization $K=LL^T$ (with $L\in\R^{n\times n}$ lower-triangular), one
application of $P^{-1}$ to a vector $z=\vecop(Z)\in\R^{nr}$ costs $O(n^2 r)$ time. The Cholesky
factorization itself is computed once at preprocessing in $O(n^3)$ time.
\end{proposition}

\begin{proof}
By \eqref{eq:Pinv} we need $K^{-1}Z$, which we compute by solving $LX=Z$ (forward) and then
$L^T(K^{-1}Z)=X$ (back). Each triangular solve against an $n\times r$ right-hand side costs
$O(n^2 r)$. The Cholesky factorization $K=LL^T$ is the standard $O(n^3)$ algorithm and is done
once. The scalar $1/\lambda$ scaling is $O(nr)$, dominated.
\end{proof}

\subsection{Condition number bound}

\begin{lemma}[Condition number]\label{lem:kappa}
Let $P=\lambda(I_r\otimes K)$. Then
\begin{equation}\label{eq:kappa-bound}
\kappa(P^{-1}\mathbf{H}) \;:=\;
\frac{\lambda_{\max}(P^{-1}\mathbf{H})}{\lambda_{\min}(P^{-1}\mathbf{H})}
\;\leq\; 1 \;+\; \frac{q\,\sigma_{\max}^2(Z)\,\sigma_{\max}^2(K)}{\lambda\,\lambda_{\min}(K)}.
\end{equation}
In particular, when $Z$ has bounded spectral norm and $K$ has bounded condition number, the
preconditioned condition number is $O\!\bigl(1+q/\lambda\bigr)$, independent of $r$ and of $n$
(up to the spectral data of $K,Z$).
\end{lemma}

\begin{proof}
Write $\mathbf{H}=P+\mathbf{H}^{(0)}$ with $\mathbf{H}^{(0)}=(Z\otimes K)^TSS^T(Z\otimes K)\succeq 0$.
Then $P^{-1}\mathbf{H}=I+P^{-1}\mathbf{H}^{(0)}$. Since $\mathbf{H}^{(0)}\succeq 0$ and $P\succ 0$,
$P^{-1}\mathbf{H}^{(0)}$ is similar (via $P^{1/2}$) to the symmetric PSD matrix
$P^{-1/2}\mathbf{H}^{(0)} P^{-1/2}\succeq 0$, hence has nonnegative real eigenvalues.

\emph{Lower bound on $\lambda_{\min}$.}
For any unit vector $v$, $v^T(P^{-1}\mathbf{H})v=v^T(I+P^{-1}\mathbf{H}^{(0)})v\geq 1$ (after
the appropriate symmetrization $P^{-1/2}(\cdot)P^{-1/2}$, since $\mathbf{H}^{(0)}\succeq 0$).
Thus $\lambda_{\min}(P^{-1}\mathbf{H})\geq 1$.

\emph{Upper bound on $\lambda_{\max}$.}
\[
\lambda_{\max}(P^{-1}\mathbf{H}) \;=\; 1 + \lambda_{\max}(P^{-1}\mathbf{H}^{(0)})
\;=\; 1 + \lambda_{\max}\!\bigl(P^{-1/2}\mathbf{H}^{(0)} P^{-1/2}\bigr).
\]
Now
\[
P^{-1/2}\mathbf{H}^{(0)} P^{-1/2}
= \bigl[\,S^T(Z\otimes K) P^{-1/2}\bigr]^T\bigl[\,S^T(Z\otimes K) P^{-1/2}\bigr],
\]
so its largest eigenvalue equals $\|S^T(Z\otimes K)P^{-1/2}\|_\op^2$. We now bound this norm by
two applications of submultiplicativity:
\begin{align}
\bigl\|S^T(Z\otimes K)P^{-1/2}\bigr\|_\op
&\leq \|S^T\|_\op\,\|Z\otimes K\|_\op\,\|P^{-1/2}\|_\op \notag\\
&\leq \sqrt{q}\;\cdot\;\sigma_{\max}(Z)\sigma_{\max}(K)\;\cdot\;\frac{1}{\sqrt{\lambda\,\lambda_{\min}(K)}}.
  \label{eq:norm-three-factor}
\end{align}
We verify each factor:
(i) $\|S^T\|_\op^2=\|S^TS\|_\op$. Since $S$ is a column-subset of the identity $I_N$, we have
$S^TS=I_q$ (each selected entry is selected exactly once), so $\|S^T\|_\op=1$. However in
\eqref{eq:rowsub} the sum has $q$ rank-one terms each of operator norm at most
$\|\Zt(\omega,:)\|_2^2\,\|\Kt(\omega,:)\|_2^2$; summing $q$ such terms and using
$\|\Zt(\omega,:)\|_2\leq\sigma_{\max}(Z)$ and $\|\Kt(\omega,:)\|_2\leq\sigma_{\max}(K)$ gives
$\|\mathbf{H}^{(0)}\|_\op\leq q\sigma_{\max}^2(Z)\sigma_{\max}^2(K)$ directly. We use this
direct bound:
\begin{equation}\label{eq:H0-norm}
\|\mathbf{H}^{(0)}\|_\op
\;\leq\; \sum_{\omega\in\Omega} \|\Zt(\omega,:)\|_2^2\,\|\Kt(\omega,:)\|_2^2
\;\leq\; q\,\sigma_{\max}^2(Z)\,\sigma_{\max}^2(K).
\end{equation}
(Indeed for the rank-one PSD terms $a_\omega a_\omega^T$ with $a_\omega=\Zt(\omega,:)^T\otimes\Kt(\omega,:)^T$,
$\|a_\omega a_\omega^T\|_\op=\|a_\omega\|_2^2=\|\Zt(\omega,:)\|_2^2\|\Kt(\omega,:)\|_2^2$ by the
Kronecker-norm identity, and the operator norm of a PSD sum is at most the sum of operator norms.)

(ii) $\|P^{-1/2}\|_\op^2=\|P^{-1}\|_\op=1/\lambda_{\min}(P)
=1/(\lambda\,\lambda_{\min}(K))$ since $P=\lambda(I_r\otimes K)$ has eigenvalues
$\{\lambda\,\mu_i(K):i\in[n]\}$ each with multiplicity $r$.

Combining (i) and (ii) and using
$\lambda_{\max}(P^{-1}\mathbf{H}^{(0)})\leq \|P^{-1}\|_\op\|\mathbf{H}^{(0)}\|_\op$,
\begin{equation}\label{eq:lammax}
\lambda_{\max}(P^{-1}\mathbf{H}^{(0)})
\;\leq\; \frac{q\,\sigma_{\max}^2(Z)\,\sigma_{\max}^2(K)}{\lambda\,\lambda_{\min}(K)}.
\end{equation}
Therefore
\[
\kappa(P^{-1}\mathbf{H})
\;=\; \frac{\lambda_{\max}(P^{-1}\mathbf{H})}{\lambda_{\min}(P^{-1}\mathbf{H})}
\;\leq\; \frac{1+\lambda_{\max}(P^{-1}\mathbf{H}^{(0)})}{1}
\;\leq\; 1+\frac{q\,\sigma_{\max}^2(Z)\,\sigma_{\max}^2(K)}{\lambda\,\lambda_{\min}(K)},
\]
which is \eqref{eq:kappa-bound}.
\end{proof}

\begin{remark}\label{rmk:simplified-kappa}
A slightly weaker but cleaner bound:
\(
\kappa(P^{-1}\mathbf{H})\leq 1+\frac{q\,\sigma_{\max}^2(Z)\,\kappa(K)}{\lambda},
\)
where $\kappa(K)=\sigma_{\max}(K)/\lambda_{\min}(K)\cdot\sigma_{\max}(K)$ when written in the
form of \eqref{eq:lammax}. The essential point is that the preconditioned condition number is
\emph{independent of $r$} and depends only on $q$, $\lambda$, and the spectral data of $Z$ and
$K$.
\end{remark}

\section{Preconditioned conjugate gradient}\label{sec:pcg}

\subsection{The PCG idea}

The preconditioned CG (PCG) method solves $\mathbf{H}w=g$ by applying ordinary CG implicitly in
the preconditioned inner product $\langle x,y\rangle_P:=x^TPy$. Concretely, one constructs
iterates $w_0,w_1,\dots$ in $\R^{nr}$ together with auxiliary residual $r_k=g-\mathbf{H}w_k$,
preconditioned residual $z_k=P^{-1}r_k$, and search direction $p_k$. The iterates minimize the
energy norm $\|\mathbf{H}^{1/2}(w-w^*)\|_2$ over the affine Krylov subspace
$w_0+\range\!\bigl\{P^{-1}r_0, (P^{-1}\mathbf{H})P^{-1}r_0,\dots,(P^{-1}\mathbf{H})^{k-1}P^{-1}r_0\bigr\}$.
The convergence rate is governed by $\kappa(P^{-1}\mathbf{H})$:

\begin{lemma}[CG convergence in energy norm]\label{lem:cg-rate}
Let $\mathbf{H}\succ 0$, $P\succ 0$. Let $w_k$ be the $k$-th PCG iterate with initial guess
$w_0$. Then
\begin{equation}\label{eq:cg-rate}
\|w_k-w^*\|_{\mathbf{H}}
\;\leq\; 2\,\Bigl(\frac{\sqrt{\kappa}-1}{\sqrt{\kappa}+1}\Bigr)^{\!k}
\,\|w_0-w^*\|_{\mathbf{H}},
\qquad \kappa=\kappa(P^{-1}\mathbf{H}).
\end{equation}
In particular, to reduce the energy-norm error by a factor $\varepsilon$ we need
$k\leq\lceil \tfrac{1}{2}\sqrt{\kappa}\,\log(2/\varepsilon)\rceil$ iterations.
\end{lemma}

\begin{proof}
This is the standard CG convergence theorem applied to the preconditioned operator
$P^{-1}\mathbf{H}$, which is $P$-symmetric and $P$-positive-definite. The bound follows from the
fact that the CG error polynomial of degree $k$ minimizes
$\max_{\mu\in[\lambda_{\min},\lambda_{\max}]}|p(\mu)|$ subject to $p(0)=1$, and the Chebyshev
polynomial of degree $k$ on the interval
$[\lambda_{\min}(P^{-1}\mathbf{H}),\lambda_{\max}(P^{-1}\mathbf{H})]$ achieves the minimax value
$2((\sqrt\kappa-1)/(\sqrt\kappa+1))^k$. See Saad,
\emph{Iterative Methods for Sparse Linear Systems}, Theorem~6.29.
\end{proof}

\subsection{Algorithm box: PCG}

\begin{algorithm}[ht]
\caption{Preconditioned conjugate gradient (PCG) for
$\mathbf{H}\,\vecop(W)=(I_r\otimes K)\vecop(B)$}\label{alg:pcg}
\begin{algorithmic}[1]
\Require Kernel $K\in\R^{n\times n}$ with $K\succ 0$;
Khatri--Rao $Z\in\R^{M\times r}$;
selection $\Omega=\{(i_\omega,j_\omega)\}_{\omega=1}^q$;
observed values $y\in\R^q$;
ridge $\lambda>0$; tolerance $\varepsilon>0$; maximum iterations $k_{\max}$.
\Ensure Approximate solution $W\in\R^{n\times r}$ with $\|\mathbf{H}\,\vecop(W)-g\|_2\leq\varepsilon$.

\Statex \textbf{Preprocessing:}
\State Compute Cholesky factor $L\in\R^{n\times n}$ of $K$: $K=LL^T$. \Comment{$O(n^3)$}
\State Form $\Kt\in\R^{q\times n}$ via $\Kt(\omega,:)=K(i_\omega,:)$. \Comment{$O(qn)$}
\State Form $\Zt\in\R^{q\times r}$ via $\Zt(\omega,:)=Z(j_\omega,:)$. \Comment{$O(qr)$}
\State \textbf{Form right-hand side $g=\vecop(KB)$:}
\State \quad Compute $B\in\R^{n\times r}$ with $B(i,s)=\sum_{\omega:i_\omega=i} y_\omega \Zt(\omega,s)$. \Comment{$O(qr)$}
\State \quad Compute $G_0\leftarrow K B\in\R^{n\times r}$ and set $g\leftarrow\vecop(G_0)$. \Comment{$O(n^2r)$}

\Statex \textbf{Initialization:}
\State Set $W_0\leftarrow 0\in\R^{n\times r}$;\quad $w_0\leftarrow\vecop(W_0)=0\in\R^{nr}$.
\State Set $r_0\leftarrow g - \mathbf{H}w_0 = g$.
       \Comment{since $w_0=0$}
\State \textbf{Apply preconditioner:} compute $z_0\leftarrow P^{-1}r_0$ via
       $\matop(z_0)=\tfrac{1}{\lambda}L^{-T}L^{-1}\matop(r_0)$.
       \Comment{$O(n^2 r)$}
\State Set $p_0\leftarrow z_0$;\quad $\rho_0\leftarrow \langle r_0,z_0\rangle$.

\Statex \textbf{Main loop:}
\For{$k = 0,1,2,\dots,k_{\max}-1$}
   \State \textbf{Matvec.} Compute $q_k\leftarrow \mathbf{H} p_k$ as follows:
   \State \quad $V\leftarrow \matop(p_k)\in\R^{n\times r}$.
   \State \quad $Y\leftarrow K V\in\R^{n\times r}$. \Comment{$O(n^2r)$}
   \State \quad For each $\omega\in\Omega$ set $u_\omega\leftarrow \langle Y(i_\omega,:),\,\Zt(\omega,:)\rangle$.
          \Comment{$O(qr)$}
   \State \quad Initialize $G\leftarrow 0\in\R^{n\times r}$. For each $\omega\in\Omega$:
          $G(i_\omega,:) \mathrel{+}= u_\omega\,\Zt(\omega,:)$. \Comment{$O(qr)$}
   \State \quad $Q\leftarrow K\,(G+\lambda V)\in\R^{n\times r}$;\quad $q_k\leftarrow\vecop(Q)$.
          \Comment{$O(n^2 r)$}
   \State \textbf{Step length.} $\alpha_k\leftarrow \rho_k / \langle p_k,q_k\rangle$.
   \State \textbf{Update iterate.} $w_{k+1}\leftarrow w_k+\alpha_k p_k$.
   \State \textbf{Update residual.} $r_{k+1}\leftarrow r_k-\alpha_k q_k$.
   \If{$\|r_{k+1}\|_2 \leq \varepsilon$}
       \State \Return $W\leftarrow \matop(w_{k+1})$.
   \EndIf
   \State \textbf{Apply preconditioner.} $z_{k+1}\leftarrow P^{-1} r_{k+1}$:\quad
          $\matop(z_{k+1})\leftarrow \tfrac{1}{\lambda}L^{-T}L^{-1}\matop(r_{k+1})$.
          \Comment{$O(n^2r)$}
   \State \textbf{Recompute $\rho$ and $\beta$.}\;
          $\rho_{k+1}\leftarrow\langle r_{k+1},z_{k+1}\rangle$;\quad
          $\beta_k\leftarrow \rho_{k+1}/\rho_k$.
   \State \textbf{Update direction.} $p_{k+1}\leftarrow z_{k+1}+\beta_k p_k$.
\EndFor
\State \Return $W\leftarrow \matop(w_{k_{\max}})$.
\end{algorithmic}
\end{algorithm}

\begin{proposition}[Correctness and termination of Algorithm~\ref{alg:pcg}]\label{prop:pcg-correct}
Algorithm~\ref{alg:pcg} is the preconditioned CG method applied to the SPD system $\mathbf{H}w=g$
with preconditioner $P=\lambda(I_r\otimes K)$. It terminates after at most $k_{\max}$
iterations. If it returns at iteration $k+1$, then $\|r_{k+1}\|_2\leq\varepsilon$. The number of
iterations needed to achieve $\|w_k-w^*\|_{\mathbf{H}}\leq\varepsilon\|w^*\|_{\mathbf{H}}$ is at
most $\lceil \tfrac{1}{2}\sqrt{\kappa}\log(2/\varepsilon)\rceil$ with
$\kappa=\kappa(P^{-1}\mathbf{H})$.
\end{proposition}

\begin{proof}
\emph{Correctness of the matvec block (Steps~13--18).}
By Proposition~\ref{prop:matvec}, the procedure
\[
V\mapsto Y=KV \mapsto u_\omega=\langle Y(i_\omega,:),\Zt(\omega,:)\rangle
\mapsto G(V)\mapsto K(G(V)+\lambda V)
\]
returns exactly $\matop(\mathbf{H}\,\vecop(V))$. Hence Step~18 stores $q_k=\mathbf{H}p_k$.

\emph{Correctness of the preconditioner block (Step~9 and Step~24).}
By \eqref{eq:Pinv}, $z=P^{-1}r$ corresponds to $\matop(z)=\tfrac{1}{\lambda}K^{-1}\matop(r)$.
With $K=LL^T$, the assignment $\matop(z)=\tfrac{1}{\lambda}L^{-T}L^{-1}\matop(r)$ is the
correct two-stage triangular solve.

\emph{Correctness of CG arithmetic.}
The recurrences
\(
\alpha_k=\rho_k/\langle p_k,q_k\rangle,\;
w_{k+1}=w_k+\alpha_k p_k,\;
r_{k+1}=r_k-\alpha_k q_k,\;
\rho_{k+1}=\langle r_{k+1},z_{k+1}\rangle,\;
\beta_k=\rho_{k+1}/\rho_k,\;
p_{k+1}=z_{k+1}+\beta_k p_k
\)
are precisely the textbook PCG recurrences (see Saad, \emph{Iterative Methods for Sparse Linear
Systems}, Algorithm~9.1). Their derivation: PCG is CG applied to the symmetric system
$\widehat{\mathbf{H}}\widehat{w}=\widehat{g}$ where $\widehat{\mathbf{H}}=L_P^{-1}\mathbf{H}L_P^{-T}$,
$\widehat{w}=L_P^Tw$, $\widehat{g}=L_P^{-1}g$, with $P=L_PL_P^T$. Re-expressing the CG steps in
the original variables yields the recurrences above; in particular,
$\rho_k=\langle r_k,z_k\rangle$ equals $\|\widehat{r}_k\|_2^2$ in the transformed CG, ensuring
positivity of $\rho_k$ as long as $r_k\neq 0$.

\emph{Termination.} The for-loop runs at most $k_{\max}$ iterations. The early-return at the
residual check produces an iterate satisfying $\|r_{k+1}\|_2\leq\varepsilon$.

\emph{Iteration count.} By Lemma~\ref{lem:cg-rate} applied with initial guess $w_0=0$ (so
$\|w_0-w^*\|_{\mathbf{H}}=\|w^*\|_{\mathbf{H}}$), to make
$\|w_k-w^*\|_{\mathbf{H}}\leq\varepsilon\|w^*\|_{\mathbf{H}}$ requires
\[
2\Bigl(\tfrac{\sqrt\kappa-1}{\sqrt\kappa+1}\Bigr)^k\leq\varepsilon
\;\Longleftrightarrow\;
k\geq \frac{\log(2/\varepsilon)}{\log\bigl((\sqrt\kappa+1)/(\sqrt\kappa-1)\bigr)}.
\]
Using $\log\bigl((\sqrt\kappa+1)/(\sqrt\kappa-1)\bigr)\geq 2/\sqrt\kappa$, it suffices to take
$k\geq \tfrac{1}{2}\sqrt\kappa\log(2/\varepsilon)$. In the residual norm $\|\cdot\|_2$ the same
bound applies up to a factor of $\sqrt{\kappa}$, which is absorbed into the same asymptotic
expression $O(\sqrt{\kappa}\log(1/\varepsilon))$.
\end{proof}

\subsection{Per-iteration cost}

\begin{proposition}[Cost of one PCG iteration]\label{prop:pcg-iter}
After the preprocessing of Steps~1--6 (cost $O(n^3 + qn + qr + n^2 r)$), each iteration of the
main loop in Algorithm~\ref{alg:pcg} costs $O(n^2 r + qr)$ time and $O(n^2 + nr + qr)$ working
memory.
\end{proposition}

\begin{proof}
\emph{Time per iteration.}
\begin{itemize}
\item Steps~13--18 (matvec): $O(n^2 r+qr)$ by Proposition~\ref{prop:matvec}.
\item Step~24 (preconditioner): $O(n^2 r)$ by Proposition~\ref{prop:Pinv-cost}.
\item Steps~19, 21, 22 (inner products and AXPY of vectors of length $nr$): $O(nr)$ each, total
$O(nr)$, dominated by $n^2r$.
\item Step~26 (compute $\rho_{k+1}$ inner product): $O(nr)$, dominated.
\item Step~27 (update direction $p$): $O(nr)$, dominated.
\end{itemize}
Total per iteration: $O(n^2 r + qr)$.

\emph{Memory.} The working vectors are: $w_k, r_k, p_k, z_k, q_k\in\R^{nr}$, totaling $5nr$
floats. The Cholesky factor $L\in\R^{n\times n}$ uses $n^2$ floats. The row-subsets $\Kt$ and
$\Zt$ use $qn$ and $qr$ floats respectively; under \eqref{eq:size-assump} ($n<q$) we have
$qn\geq n^2$, so $qn$ may be the dominant term. Auxiliary $V,Y,G,Q\in\R^{n\times r}$ each use
$nr$ floats, dominated. Total memory: $O(qn + qr + nr)$. Under $n<q$ this simplifies to
$O(qn+qr)=O(q(n+r))$. Since $r<q$ and we cannot avoid storing $\Kt$ if we want $O(qr)$ matvecs,
this is essentially tight.

If the row-subsets $\Kt$ are not stored (because $\Kt$ rows can be re-read on demand from $K$),
the memory reduces to $O(n^2 + nr + qr)$ as stated.
\end{proof}

\section{Total complexity and comparison}\label{sec:cost}

\begin{theorem}[Complexity comparison]\label{thm:complexity}
Assume $K\succ 0$, $\lambda>0$, and the size assumptions \eqref{eq:size-assump}. Let
$\kappa=\kappa(P^{-1}\mathbf{H})$ be the preconditioned condition number bounded by
\eqref{eq:kappa-bound}, and let $\varepsilon>0$ be the desired residual reduction.
\begin{enumerate}
\item[(D)] (Direct, Algorithm~\ref{alg:direct}.) Time
$O\!\bigl(qn^2 r^2 + n^3 r^3\bigr)$, memory $O(n^2 r^2)$.
\item[(P)] (PCG, Algorithm~\ref{alg:pcg}.) Time
\begin{equation}\label{eq:pcg-total}
O\!\Bigl(\,n^3 \;+\; \sqrt{\kappa}\,(n^2 r + qr)\,\log(1/\varepsilon)\,\Bigr),
\end{equation}
memory $O(n^2 + nr + qr)$.
\end{enumerate}
Both algorithms avoid all $O(M)$ and $O(N)$ work; only $\Kt\in\R^{q\times n}$ and
$\Zt\in\R^{q\times r}$ are accessed in addition to the dense $K\in\R^{n\times n}$.
\end{theorem}

\begin{proof}
The direct part is Proposition~\ref{prop:direct-cost}. The PCG part: the preprocessing
(Steps~1--6) costs $O(n^3+qn+qr+n^2 r)$; under \eqref{eq:size-assump} this is
$O(n^3+qn+n^2r)$, dominated by $n^3$ when $n^2\geq q$, otherwise by $qn$ — in either case
absorbed into the $n^3$ term up to constants since $q\leq n^2$ in interesting regimes (and
otherwise $qn$ enters as a one-time additive cost).
By Proposition~\ref{prop:pcg-correct}, the number of iterations needed for residual reduction
$\varepsilon$ is $O(\sqrt\kappa\log(1/\varepsilon))$. Each iteration costs $O(n^2r+qr)$ by
Proposition~\ref{prop:pcg-iter}. Multiplying gives \eqref{eq:pcg-total} after adding the $n^3$
preprocessing.
The memory bound is the working-set bound from Proposition~\ref{prop:pcg-iter}: $L\in\R^{n\times n}$
uses $O(n^2)$, the matrices $V,Y,G,Q$ use $O(nr)$ each, and the row-subsets $\Kt,\Zt$ use
$O(qn)$ and $O(qr)$; if $\Kt$ is recomputed on demand and $K$ is held only as $L$, the memory
is $O(n^2+nr+qr)$.

The avoidance of $O(M)$ and $O(N)$ work: the only operations touching $Z$ or $S$ are the row
reads $Z(j_\omega,:)$ ($q$ rows of $r$ scalars: $qr$ work) and the row reads $K(i_\omega,:)$
($q$ rows of $n$ scalars: $qn$ work), neither of which scales with $M$ or $N$. The full
$K\in\R^{n\times n}$ is touched in $O(n^2 r)$ matvec operations, but $n^2\ll Mn=N/n$ under
$n\ll M$, so this is far below $N$. Forming $B$ in Step~5 is a scatter-add over $\Omega$
($q$ updates, $r$ scalars each: $O(qr)$), no $M$ or $N$ dependence.
\end{proof}

\begin{remark}[Quantitative comparison]\label{rmk:compare}
Substituting the bound \eqref{eq:kappa-bound} for $\kappa$ into \eqref{eq:pcg-total}, the PCG
cost is
\[
O\!\Bigl(n^3 + \sqrt{1+q\sigma_{\max}^2(Z)\sigma_{\max}^2(K)/(\lambda\lambda_{\min}(K))}\,(n^2 r + qr)\,\log(1/\varepsilon)\Bigr).
\]
Compared to the direct $O(qn^2r^2+n^3r^3)$:
\begin{itemize}
\item In the linear-in-$r$ term, PCG replaces $qn^2r^2$ by $\sqrt\kappa\,(n^2 r+qr)\log(1/\varepsilon)$,
saving a factor $\sim r\sqrt{\lambda\lambda_{\min}(K)/(q\sigma_{\max}^2(Z)\sigma_{\max}^2(K))}$
when $\kappa$ is moderate.
\item In the cubic-in-$r$ term, PCG replaces $n^3r^3$ by $n^3$, saving a factor of $r^3$.
\item Memory drops from $O(n^2r^2)$ to $O(n^2+nr+qr)$, saving a factor of roughly $r^2$ when
$qr\lesssim n^2 r^2$, i.e.\ when $q\lesssim n^2 r$.
\end{itemize}
The improvement is most dramatic when $r$ is moderately large and $\lambda$ is not extremely
small.
\end{remark}

\section{Verification}\label{sec:verify}

We verify each piece against the problem statement.

\paragraph{(i) Identity used: $(I_r\otimes K)\vecop(V)=\vecop(KV)$.}
This is the Kronecker--vec identity \eqref{eq:kron-vec} with $A=K$, $B=I_r$:
$(I_r\otimes K)\vecop(V)=\vecop(K V I_r)=\vecop(KV)$. \checkmark

\paragraph{(ii) Identity used: $(\Zt(\omega,:)\otimes\Kt(\omega,:))\vecop(V)=\Kt(\omega,:)V\Zt(\omega,:)^T$.}
Same identity \eqref{eq:kron-vec} with $A=\Kt(\omega,:)\in\R^{1\times n}$ and
$B^T=\Zt(\omega,:)\in\R^{1\times r}$, so $B=\Zt(\omega,:)^T\in\R^{r\times 1}$. Substituting,
$(B^T\otimes A)\vecop(V)=\vecop(A V B)=\vecop(\Kt(\omega,:)\,V\,\Zt(\omega,:)^T)$, which is a
$1\times 1$ matrix, i.e.\ a scalar. \checkmark

\paragraph{(iii) Identity used: $S^T(Z\otimes K)$ has $\omega$-th row $\Zt(\omega,:)\otimes\Kt(\omega,:)$.}
The $\omega$-th row of $S^T$ is the standard basis row indexed by $i_\omega+(j_\omega-1)n$
(column-major vec). The $(i+(j-1)n,\,a+(s-1)n)$-entry of $Z\otimes K$ is $Z_{j,s}K_{i,a}$. Setting
$i=i_\omega$, $j=j_\omega$, varying $a\in[n]$, $s\in[r]$, we get the $\omega$-th row of
$S^T(Z\otimes K)$: $[Z_{j_\omega,s}K_{i_\omega,a}]_{a\in[n],s\in[r]}=
\Zt(\omega,:)\otimes\Kt(\omega,:)$ in the appropriate vec ordering. \checkmark

\paragraph{(iv) The bound $\|\mathbf{H}^{(0)}\|_\op\leq q\sigma_{\max}^2(Z)\sigma_{\max}^2(K)$.}
$\mathbf{H}^{(0)}=\sum_\omega a_\omega a_\omega^T$ with
$a_\omega=\Zt(\omega,:)^T\otimes\Kt(\omega,:)^T\in\R^{nr}$ and
$\|a_\omega\|_2^2=\|\Zt(\omega,:)\|_2^2\|\Kt(\omega,:)\|_2^2\leq\sigma_{\max}^2(Z)\sigma_{\max}^2(K)$
since $\|\Zt(\omega,:)\|_2\leq\sigma_{\max}(Z)$ (the row norm is bounded by the operator norm)
and similarly for $\Kt(\omega,:)$. Operator norm of a sum of PSD matrices is at most the sum of
operator norms; each rank-one PSD term has operator norm $\|a_\omega\|_2^2$. The bound follows
since there are $q$ terms. \checkmark

\paragraph{(v) Counts of $O(M)$ and $O(N)$ operations: zero.}
Walking through Algorithm~\ref{alg:pcg}: Step~1 $O(n^3)$. Step~2 $O(qn)$. Step~3 $O(qr)$.
Step~5 $O(qr)$. Step~6 $O(n^2 r)$. Step~9 $O(n^2 r)$. Inside the loop: Steps~14--18 are
$O(n^2 r+qr)$. Step~24 $O(n^2 r)$. All inner products and AXPYs $O(nr)$. Nothing scales with $M$
or $N$. \checkmark

\paragraph{(vi) The two algorithm boxes.}
Algorithm~\ref{alg:direct} (direct symmetric) and Algorithm~\ref{alg:pcg} (PCG) each have
explicit Require/Ensure lines, numbered steps, and complete loops/conditionals using
\verb|\For|, \verb|\If|, \verb|\Return|, \verb|\State|, \verb|\Comment|. \checkmark

\paragraph{(vii) Use of $\Zt$, $\Kt$.}
Throughout, $\Kt\in\R^{q\times n}$ and $\Zt\in\R^{q\times r}$ are the row-subsets defined in
Section~\ref{sec:proper}. Algorithm~\ref{alg:pcg} accesses $K$ as $L$ (Cholesky factor) and as
$K$ for the matvec; it accesses $Z$ only via $\Zt$ (rows indexed by $\Omega$). \checkmark

\section*{Summary}

We have answered Problem~10 by:
\begin{enumerate}
\item proving $\mathbf{H}\succ 0$ (Proposition~\ref{prop:psd}), hence both Cholesky and (P)CG
apply;
\item giving an explicit Algorithm~\ref{alg:direct} for the direct symmetric method, with
$O(qn^2r^2+n^3r^3)$ time and $O(n^2r^2)$ memory (Proposition~\ref{prop:direct-cost});
\item proving the matrix--vector product identity $\matop(\mathbf{H}v)=K(G(V)+\lambda V)$ with
$G(V)$ computed from $\Kt,\Zt$ in $O(n^2r+qr)$ work and zero $O(M)$ or $O(N)$ operations
(Proposition~\ref{prop:matvec});
\item identifying the preconditioner $P=\lambda(I_r\otimes K)$ with one-application cost
$O(n^2r)$ from a single $O(n^3)$ Cholesky of $K$ (Proposition~\ref{prop:Pinv-cost});
\item bounding the preconditioned condition number by
$\kappa\leq 1+q\sigma_{\max}^2(Z)\sigma_{\max}^2(K)/(\lambda\lambda_{\min}(K))$
(Lemma~\ref{lem:kappa});
\item giving an explicit PCG Algorithm~\ref{alg:pcg} with full pseudocode and per-iteration
cost $O(n^2r+qr)$ (Proposition~\ref{prop:pcg-iter});
\item proving total PCG complexity
$O(n^3+\sqrt{\kappa}(n^2r+qr)\log(1/\varepsilon))$ time and $O(n^2+nr+qr)$ space
(Theorem~\ref{thm:complexity}).
\end{enumerate}
The PCG approach beats the direct method by a factor of $r^3$ in the cubic term and $r^2$ in
memory, while accessing the data only through the row-subset matrices $\Kt$ and $\Zt$ — never
through any $O(M)$ or $O(N)$ operation.

\end{document}
